% 王豪 — 中文简历
% 黑白打印版；字体排版与 https://wanghao9610.github.io/ 保持一致（衬线标题）。
% 编译：make -f Makefile.zh   (latexmk -xelatex)

\documentclass[UTF8,10pt,letterpaper,fontset=fandol,linespread=1.1]{ctexart}

\usepackage[margin=0.6in,top=0.5in,bottom=0.5in]{geometry}
\setmainfont{SourceSerifPro-Regular.otf}[
  BoldFont=SourceSerifPro-Semibold.otf,
  ItalicFont=SourceSerifPro-RegularIt.otf,
  BoldItalicFont=SourceSerifPro-SemiboldIt.otf]
\setsansfont{SourceSansPro-Regular.otf}[
  BoldFont=SourceSansPro-Semibold.otf,
  ItalicFont=SourceSansPro-RegularIt.otf,
  BoldItalicFont=SourceSansPro-SemiboldIt.otf]
\usepackage{xcolor}
\usepackage{hyperref}
\usepackage{enumitem}
\usepackage{tabularx}
\usepackage{titlesec}
\usepackage{array}
\usepackage{ragged2e}
\usepackage{microtype}
\usepackage{needspace}
\usepackage{fontawesome5}
\usepackage{etoolbox}

% ---- 黑白打印配色 ------------------------------------------------------------
\definecolor{ink}{HTML}{000000}
\definecolor{text}{HTML}{000000}
\definecolor{muted}{HTML}{3A3A3A}
\definecolor{soft}{HTML}{555555}
\definecolor{accent}{HTML}{000000}
\definecolor{rule}{HTML}{BFBFBF}

\hypersetup{
  colorlinks=true,
  urlcolor=accent,
  linkcolor=accent,
  pdftitle={王豪 - 简历},
  pdfauthor={王豪}
}

\pagestyle{empty}
\color{text}
\setlength{\parindent}{0pt}
\setlength{\parskip}{0pt}
\setlist[itemize]{leftmargin=1.15em,itemsep=1.4pt,topsep=2pt,parsep=0pt,partopsep=0pt,label=\textcolor{accent}{\scriptsize\raisebox{.25ex}{$\bullet$}}}
\renewcommand{\arraystretch}{1.06}
\setlength{\emergencystretch}{1.5em}

\titleformat{\section}
  {\rmfamily\Large\bfseries\color{ink}}
  {}
  {0pt}
  {}
  [{\vspace{-3pt}\color{rule}\titlerule[0.6pt]}]
\titlespacing*{\section}{0pt}{9pt}{5pt}

\newcommand{\datebadge}[1]{\textcolor{accent}{\footnotesize\bfseries #1}}
\newcommand{\metaline}[1]{\textcolor{muted}{\small #1}}
\newcommand{\venue}[1]{\textcolor{accent}{\footnotesize\bfseries #1}}
\newcommand{\stars}[1]{\ifstrempty{#1}{}{\hfill\textcolor{soft}{\footnotesize\faStar\,#1}}}
\newcommand{\sep}{{\color{rule}\;\textbar\;}}

\newcommand{\cvitem}[4]{%
  \Needspace{6\baselineskip}
  \vspace{3pt}
  \begin{tabularx}{\textwidth}{@{}>{\RaggedRight\arraybackslash}X>{\RaggedLeft\arraybackslash}p{0.25\textwidth}@{}}
    \textbf{\textcolor{ink}{#1}} & \datebadge{#2}
  \end{tabularx}
  \vspace{-1pt}
  \metaline{#3}\par
  \vspace{1pt}
  #4
}

\newcommand{\publication}[5]{%
  \Needspace{5\baselineskip}
  \vspace{2pt}
  \begin{tabularx}{\textwidth}{@{}>{\RaggedRight\arraybackslash}X>{\RaggedLeft\arraybackslash}p{0.16\textwidth}@{}}
    \textbf{\href{#1}{\textcolor{ink}{#2}}} & \venue{#3}
  \end{tabularx}
  \vspace{-1pt}
  {\small #4}\par
  \vspace{1pt}
  \textcolor{muted}{\small #5}\par
  \vspace{4.5pt}
}

\newcommand{\repo}[5]{%
  \textbf{\href{#1}{\textcolor{ink}{#2}}}\ \venue{#3}\stars{#4}\par
  \vspace{-1pt}
  \textcolor{muted}{\small #5}\par
  \vspace{3pt}
}

\begin{document}

% ============================ 抬头 ==========================================
\begin{minipage}{\textwidth}
  \vspace{2pt}
  \centering
  {\rmfamily\fontsize{28}{32}\selectfont\bfseries\color{ink} 王豪\hspace{0.35em}%
   {\fontsize{20}{24}\selectfont\mdseries\color{muted} Hao Wang}}\par
  \vspace{5pt}
  {\rmfamily\itshape\large\color{muted} Leave the world a little different.}\par
  \vspace{4pt}
  {\normalsize\color{accent} 博士研究生 · 计算机视觉 · 多模态人工智能 · 智能体}\par
  \vspace{7pt}
  {\small\color{muted}
    \faEnvelope\ \href{mailto:wanghao9610@gmail.com}{wanghao9610@gmail.com}
    \sep
    \faGlobe\ \href{https://wanghao9610.github.io/}{wanghao9610.github.io}
    \sep
    \faGithub\ \href{https://github.com/wanghao9610}{wanghao9610}
    \sep
    \faWeixin\ wangh9610
    \sep
    \faMapMarker*\ 深圳
  }\par
  \vspace{7pt}
  {\color{ink}\rule{\linewidth}{0.8pt}}\par
  \vspace{5pt}
  \begin{minipage}[t]{0.62\textwidth}
    \RaggedRight\small\color{muted} 中山大学 HCP 实验室 \& 鹏城实验室
  \end{minipage}%
  \hfill
  \begin{minipage}[t]{0.36\textwidth}
    \RaggedLeft\small\color{muted} 预计毕业时间：\textbf{\color{accent}2026 年 12 月}
  \end{minipage}\par
  \vspace{1pt}
  {\small\color{muted}\RaggedRight 寻求工业界研究岗位，同时欢迎科研合作。\par}
\end{minipage}

% ============================ 简介 ==========================================
\section{个人简介}
博士研究生，研究聚焦\textbf{开放场景计算机视觉}与\textbf{多模态大语言模型}，并持续探索\textbf{多模态智能体模型}。一作工作包括 X2SAM（ECCV 2026）、X-SAM（AAAI 2026）、OV-DINO 与 TMANet（ICIP 2021），覆盖图像/视频统一分割、开放词汇检测与时序一致的像素级理解。代表性工作均开源并配有项目主页（GitHub 累计 1.1k+ stars）。关注如何将基础模型与可靠的密集感知、智能体工作流结合，服务于实际视觉系统。

\section{研究亮点}
\begin{itemize}
  \item 建立以第一作者为主的开放场景感知研究线：从视频语义分割到开放词汇检测，再到多模态任意分割。
  \item 构建同时支持自然语言指令与视觉提示的多模态分割大模型，连接高层语义推理与密集掩码预测。
  \item 借助掩码记忆机制将任意分割从图像扩展到视频，实现跨帧时序一致的对象级与像素级理解。
  \item 长期维护开源研究代码及面向智能体的研究工具链（STAR / STORY / STAGE），重视可复现与开放共享。
\end{itemize}

% ============================ 论文 ==========================================
\section{一作论文}
\publication
  {https://wanghao9610.github.io/X2SAM/}
  {X2SAM: Any Segmentation in Images and Videos}
  {ECCV 2026}
  {\textbf{Hao Wang}, Limeng Qiao, Chi Zhang, Guanglu Wan, Lin Ma, Xiangyuan Lan, Xiaodan Liang}
  {面向图像与视频的统一分割多模态大模型，支持对话式指令与视觉提示，通过掩码记忆实现时序一致的像素级感知，将任意分割从静态图像扩展到视频。\href{https://wanghao9610.github.io/X2SAM/}{主页} \sep \href{https://arxiv.org/abs/2605.00891}{论文} \sep \href{https://github.com/wanghao9610/X2SAM}{代码}}

\publication
  {https://wanghao9610.github.io/X-SAM/}
  {X-SAM: From Segment Anything to Any Segmentation}
  {AAAI 2026}
  {\textbf{Hao Wang}, Limeng Qiao, Zequn Jie, Zhijian Huang, Chengjian Feng, Qingfang Zheng, Lin Ma, Xiangyuan Lan, Xiaodan Liang}
  {统一的多模态大模型框架，将 Segment Anything 的能力扩展到任意分割与像素级感知理解，支持指令驱动的交互式分割。\href{https://wanghao9610.github.io/X-SAM/}{主页} \sep \href{https://arxiv.org/abs/2508.04655}{论文} \sep \href{https://github.com/wanghao9610/X-SAM}{代码}}

\publication
  {https://wanghao9610.github.io/OV-DINO/}
  {OV-DINO: Unified Open-Vocabulary Detection with Language-Aware Selective Fusion}
  {arXiv 2024}
  {\textbf{Hao Wang}, Pengzhen Ren, Zequn Jie, Xiao Dong, Chengjian Feng, Yinlong Qian, Lin Ma, Dongmei Jiang, Yaowei Wang, Xiangyuan Lan, Xiaodan Liang}
  {在多源大规模数据上预训练的统一开放词汇检测器，提出语言感知的选择性融合机制，通过视觉--语言对齐提升类别泛化能力。\href{https://wanghao9610.github.io/OV-DINO/}{主页} \sep \href{https://arxiv.org/abs/2407.07844}{论文} \sep \href{https://github.com/wanghao9610/OV-DINO}{代码}}

\publication
  {https://arxiv.org/abs/2102.08643}
  {TMANet: Temporal Memory Attention for Video Semantic Segmentation}
  {ICIP 2021}
  {\textbf{Hao Wang}, Weining Wang, Jing Liu}
  {面向长程视频语义分割的时序记忆注意力方法，无需光流预测即可建模跨帧关系。\href{https://arxiv.org/abs/2102.08643}{论文} \sep \href{https://github.com/wanghao9610/TMANet}{代码}}

% ============================ 开源 ==========================================
\Needspace*{16\baselineskip}
\section{开源项目}
\metaline{研究代码与工具均托管于 \href{https://github.com/wanghao9610}{GitHub}；star 数截至 2026 年 9 月。}\par
\vspace{4pt}
\begin{minipage}[t]{0.485\textwidth}
  \repo{https://github.com/wanghao9610/OV-DINO}{OV-DINO}{arXiv 2024}{414}{语言感知选择性融合的统一开放词汇检测。}
  \repo{https://github.com/wanghao9610/X-SAM}{X-SAM}{AAAI 2026}{391}{从 Segment Anything 到任意分割。}
  \repo{https://github.com/wanghao9610/TMANet}{TMANet}{ICIP 2021}{128}{面向视频语义分割的时序记忆注意力。}
  \repo{https://github.com/wanghao9610/X2SAM}{X2SAM}{ECCV 2026}{104}{单一多模态大模型实现图像与视频任意分割。}
\end{minipage}%
\hfill
\begin{minipage}[t]{0.485\textwidth}
  \repo{https://github.com/wanghao9610/STAR}{STAR}{Harness · WIP}{51}{Systematic Toolchain for AI Research：面向 AI 研究的智能体工具链，覆盖规划、执行与追踪。}
  \repo{https://github.com/wanghao9610/STORY}{STORY \,/\, \href{https://github.com/wanghao9610/STAGE}{STAGE}}{Harness · WIP}{}{配套工具链：跨年度组织研究，以及写作、引导与编辑。}
  \repo{https://omdsh-plugins.github.io/}{omdsh-plugins}{Collection}{}{Oh My DeepSeek Harness：DeepSeek Harness 插件精选集。}
  \repo{https://github.com/wanghao9610/arXivTeX}{arXivTeX}{Template}{33}{arXiv / 会议论文模板。另有 \href{https://github.com/wanghao9610/Claude-Model-Proxy}{Claude Model Proxy}、\href{https://github.com/wanghao9610/Codex-Remote-Connector}{Codex Remote Connector}。}
\end{minipage}

% ============================ 经历 ==========================================
\section{研究经历}
\cvitem
  {美团 · M17-MM}
  {2025.01 -- 至今}
  {研究实习生}
  {\begin{itemize}
    \item 主导 X2SAM：面向图像与视频的统一任意分割多模态大模型，支持对话式指令与视觉提示。
    \item 设计掩码记忆机制，提升视频分割与交互式感知中的对象级时序一致性；负责训练、评测与开源发布。
  \end{itemize}}

\cvitem
  {美团 · 视觉智能部}
  {2022.07 -- 2025.01}
  {研究实习生}
  {\begin{itemize}
    \item 主导 OV-DINO 与 X-SAM，推进开放词汇识别与任意分割能力。
    \item 研究语言感知融合与多模态监督，用于可扩展的视觉识别与分割。
  \end{itemize}}

\cvitem
  {腾讯 · AI 平台部}
  {2021.05 -- 2021.08}
  {应用研究实习生}
  {}

\cvitem
  {华为 · 照片处理部}
  {2019.09 -- 2020.07}
  {应用项目实习生}
  {}

% ============================ 教育 ==========================================
\section{教育经历}
\cvitem
  {中山大学 \& 鹏城实验室}
  {2022.09 -- 至今}
  {博士研究生，智能工程学院}
  {\begin{itemize}
    \item 导师：梁小丹教授、蓝湘源副教授。
    \item 研究方向：图像/视频分割、开放场景视觉感知、多模态基础模型与多模态智能体模型；预计 2026 年 12 月毕业。
  \end{itemize}}

\cvitem
  {中国科学院大学 \& 中国科学院自动化研究所}
  {2019.09 -- 2022.06}
  {硕士研究生，人工智能学院；导师：刘静教授}
  {}

\cvitem
  {北京交通大学}
  {2015.09 -- 2019.06}
  {本科生，电子信息工程学院}
  {}

% ============================ 其他 ==========================================
\section{合作论文}
\publication
  {https://ieeexplore.ieee.org/abstract/document/10096063}
  {WL-MSR: Watch and Listen for Multimodal Subtitle Recognition}
  {ICASSP 2023}
  {Jiawei Liu, \textbf{Hao Wang}, Weining Wang, Xingjian He, Jing Liu}
  {基于 Transformer 的多模态字幕识别方法，融合 OCR 与 ASR 信息，并使用掩码/裁剪策略和多层级身份嵌入。}

\Needspace*{5\baselineskip}
\section{学术服务与获奖}
\begin{tabular}{@{}>{\bfseries\color{ink}}l@{\hspace{0.08\textwidth}}l@{}}
会议审稿人 & AAAI 2026、ECCV 2024、ICCV 2023 \\
期刊审稿人 & Proceedings of the IEEE \\
获奖 & ICCV 2021 第 1 届 VSPW Challenge Workshop 第一名（2021.09）
\end{tabular}

\end{document}
